\documentclass{article}
\usepackage{amsmath,amssymb,amsthm,algorithm,algorithmic}
\usepackage{hyperref}
\usepackage{enumitem}
\newtheorem{theorem}{Theorem}
\newtheorem{lemma}{Lemma}
\newtheorem{assumption}{Assumption}
\newtheorem{remark}{Remark}
\begin{document}

\section*{AMSGrad (Adam with Monotone Second‑Moment Estimates)}

\section*{Mathematical Formulation}
Let $f:\mathbb{R}^d\rightarrow\mathbb{R}$ be the objective function.  
At iteration $t\ge 1$ we have a stochastic gradient $g_t:=\nabla f(w_t;\xi_t)$ where $\xi_t$ is a random sample.  
Hyper‑parameters: learning rate $\eta>0$, exponential decay rates $0<\beta_1,\beta_2<1$, and a small constant $\varepsilon>0$.

\[
\begin{aligned}
m_t &= \beta_1 m_{t-1} + (1-\beta_1) g_t, \qquad &&\text{(first moment)}\\[4pt]
v_t &= \beta_2 v_{t-1} + (1-\beta_2) g_t\odot g_t, \qquad &&\text{(second moment)}\\[4pt]
\hat{v}_t &= \max\!\big(\hat{v}_{t-1},\, v_t\big) \quad\text{(elementwise max)}\\[4pt]
\hat{m}_t &= \frac{m_t}{1-\beta_1^{t}}\quad\text{(bias‑corrected first moment)}\\[4pt]
w_{t+1} &= w_t - \eta\,\frac{\hat{m}_t}{\sqrt{\hat{v}_t}+\varepsilon},
\end{aligned}
\]
with the convention $m_0=v_0=\hat{v}_0=0$.

\section*{Pseudocode}
\begin{algorithm}[H]
\caption{AMSGrad}
\begin{algorithmic}[1]
\REQUIRE Initial point $w_0\in\mathbb{R}^d$, stepsize $\eta$, decay rates $\beta_1,\beta_2$, tolerance $\varepsilon$.
\STATE $m_0\gets 0,\; v_0\gets 0,\; \hat{v}_0\gets 0$
\FOR{$t=1,2,\dots,T$}
    \STATE Sample a mini‑batch (or a data point) and compute stochastic gradient $g_t=\nabla f(w_t;\xi_t)$
    \STATE $m_t \gets \beta_1 m_{t-1} + (1-\beta_1) g_t$
    \STATE $v_t \gets \beta_2 v_{t-1} + (1-\beta_2) g_t\odot g_t$
    \STATE $\hat{v}_t \gets \max(\hat{v}_{t-1},\,v_t)$ \hfill\textit{(elementwise)}
    \STATE $\hat{m}_t \gets m_t/(1-\beta_1^{t})$
    \STATE $w_{t+1} \gets w_t - \eta\,\hat{m}_t/(\sqrt{\hat{v}_t}+\varepsilon)$
\ENDFOR
\RETURN $w_{T+1}$
\end{algorithmic}
\end{algorithm}

\section*{Dry Run Example}
Consider the scalar quadratic $f(w)=(w-3)^2$, thus $\nabla f(w)=2(w-3)$.  
Take $w_0=0$, $\eta=0.1$, $\beta_1=0.9$, $\beta_2=0.999$, $\varepsilon=10^{-8}$ and assume **exact** gradients (so $g_t=\nabla f(w_t)$).

\begin{center}
\begin{tabular}{c|c|c|c|c|c}
$t$ & $w_t$ & $g_t$ & $m_t$ & $v_t$ & $w_{t+1}$ \\ \hline
0 & $0$ & – & $0$ & $0$ & – \\
1 & $0$ & $-6$ & $0.1\cdot(-6)=-0.6$ & $0.001\cdot36=0.036$ & $0-0.1\frac{-0.6}{\sqrt{0.036}}=0+0.1\cdot1=0.1$ \\
2 & $0.1$ & $-5.8$ & $0.9(-0.6)+0.1(-5.8)=-1.14$ & $0.999\cdot0.036+0.001\cdot33.64=0.0696$ & $\hat v_2=\max(0.036,0.0696)=0.0696$\\
  &  &  &  &  & $w_2=0.1-0.1\frac{-1.14}{\sqrt{0.0696}}=0.1+0.1\cdot4.33\approx0.533$\\
3 & $0.533$ & $-4.934$ & $0.9(-1.14)+0.1(-4.934)=-1.5306$ & $0.999\cdot0.0696+0.001\cdot24.34=0.0942$ & $\hat v_3=\max(0.0696,0.0942)=0.0942$\\
  &  &  &  &  & $w_3=0.533-0.1\frac{-1.5306}{\sqrt{0.0942}}=0.533+0.1\cdot4.99\approx1.032$
\end{tabular}
\end{center}
The objective values are $f(w_1)=9$, $f(w_2)\approx6.17$, $f(w_3)\approx3.94$, showing monotone decrease.

\newpage
\section*{Assumptions}
\begin{assumption}[Smoothness]\label{ass:smooth}
$f$ is $L$‑smooth, i.e. for all $x,y\in\mathbb{R}^d$,
\[
\|\nabla f(x)-\nabla f(y)\|\le L\|x-y\|.
\]
Equivalently,
\[
f(y)\le f(x)+\langle\nabla f(x),y-x\rangle+\frac{L}{2}\|y-x\|^2.
\]
\end{assumption}
\begin{assumption}[Bounded Stochastic Gradients]\label{ass:bound}
There exists $G>0$ such that
\[
\mathbb{E}_{\xi_t}\big[\|g_t\|^2\big]\le G^2,\qquad\forall t.
\]
\end{assumption}
\begin{assumption}[Unbiasedness]\label{ass:unbiased}
\[
\mathbb{E}_{\xi_t}[g_t\mid \mathcal{F}_{t-1}]=\nabla f(w_t),
\]
where $\mathcal{F}_{t-1}$ is the $\sigma$‑algebra generated by $\{w_1,\dots,w_{t-1}\}$.
\end{assumption}
\begin{assumption}[Lower Boundedness]\label{ass:lb}
$f$ has a finite infimum $f^\star:=\inf_{w}f(w)>-\infty$.
\end{assumption}

\section*{Main Convergence Theorem}
\begin{theorem}[Convergence of AMSGrad for non‑convex $L$‑smooth objectives]\label{thm:main}
Assume \ref{ass:smooth}–\ref{ass:lb} and let the stepsize be chosen as $\eta_t=\frac{\eta}{\sqrt{t}}$ with a constant $\eta>0$.  
Define $d$ the dimension and $\beta_1,\beta_2$ satisfy $0<\beta_1<\beta_2<1$ and $\frac{\beta_1}{\sqrt{\beta_2}}<1$.  
Then for any $T\ge 1$,
\[
\frac{1}{T}\sum_{t=1}^{T}\mathbb{E}\big[\|\nabla f(w_t)\|^2\big]
\;\le\;
\frac{C_1}{\sqrt{T}}+\frac{C_2\log T}{\sqrt{T}},
\]
where
\[
C_1=\frac{2\big(f(w_1)-f^\star\big)}{\eta(1-\beta_1)}\sqrt{d},\qquad
C_2=\frac{L\eta G^2}{(1-\beta_1)(1-\beta_2)}\sqrt{d}.
\]
Consequently $\displaystyle\min_{1\le t\le T}\mathbb{E}\big[\|\nabla f(w_t)\|^2\big]=O\!\big(\tfrac{\log T}{\sqrt{T}}\big)$.
\end{theorem}

\section*{Theoretical Properties}
\begin{itemize}
\item \textbf{Stability via monotone $\hat v_t$:} the max operation guarantees $\hat v_t\ge\hat v_{t-1}$ elementwise, preventing the denominator from shrinking and thus avoiding excessively large steps.
\item \textbf{Hyper‑parameter sensitivity:} the bound depends on $\frac{1}{1-\beta_1}$ and $\frac{1}{1-\beta_2}$; choosing $\beta_1$ close to $1$ (as common in practice) inflates $C_1$ but does not affect the $O(1/\sqrt{T})$ rate.
\item \textbf{Comparison to Adam:} Adam without the max operation can diverge on simple convex examples (Reddi et al., 2018). AMSGrad restores a provable $O(1/\sqrt{T})$ rate while preserving Adam’s adaptive per‑coordinate stepsizes.
\item \textbf{Extension to mini‑batch gradients:} the same proof holds with $G^2$ set to the bound on the second moment of the mini‑batch gradient.
\end{itemize}

\section*{Discussion}
The theorem shows that despite the non‑convex landscape, AMSGrad drives the average squared gradient norm to zero at the same rate as classical stochastic gradient methods with diminishing stepsizes. The monotone second‑moment estimate is the key device that eliminates the “exploding stepsize’’ pathology identified for Adam.

\newpage
\section*{Preliminaries (Definitions and Lemmas)}
\begin{lemma}[Smoothness inequality]\label{lem:smooth}
For an $L$‑smooth function,
\[
f(w_{t+1})\le f(w_t)+\langle\nabla f(w_t),w_{t+1}-w_t\rangle+\frac{L}{2}\|w_{t+1}-w_t\|^2 .
\]
\end{lemma}
\begin{proof}
Directly from the definition of $L$‑smoothness. \hfill $\square$
\end{proof}

\begin{lemma}[Bound on the adaptive denominator]\label{lem:denom}
For all $t$ and each coordinate $i$,
\[
\sqrt{\hat v_{t,i}} \ge \sqrt{(1-\beta_2)v_{1,i}}.
\]
\end{lemma}
\begin{proof}
Since $v_{t,i}=\beta_2 v_{t-1,i}+(1-\beta_2)g_{t,i}^2\ge (1-\beta_2)g_{t,i}^2$ and $\hat v_t\ge v_t$ elementwise, the claim follows by induction. \hfill $\square$
\end{proof}

\section*{Main Proof (Step‑by‑Step Derivation)}
We prove Theorem~\ref{thm:main}. All expectations are taken w.r.t. the randomness up to iteration $t$ unless otherwise noted.

\noindent\textbf{Step 1 (Apply smoothness).}  
From Lemma~\ref{lem:smooth} with $w_{t+1}=w_t-\eta_t\frac{\hat m_t}{\sqrt{\hat v_t}+\varepsilon}$,
\begin{align}
f(w_{t+1}) &\le f(w_t)-\eta_t\Big\langle\nabla f(w_t),\frac{\hat m_t}{\sqrt{\hat v_t}+\varepsilon}\Big\rangle
            +\frac{L\eta_t^2}{2}\Big\|\frac{\hat m_t}{\sqrt{\hat v_t}+\varepsilon}\Big\|^2 .
\tag{1}\label{eq:smooth}
\end{align}
\textbf{[Step 1]}

\noindent\textbf{Step 2 (Decompose $\hat m_t$).}  
Recall $\hat m_t=m_t/(1-\beta_1^{t})$ and $m_t$ is a convex combination of past stochastic gradients:
\[
m_t = (1-\beta_1)\sum_{i=1}^{t}\beta_1^{t-i} g_i .
\]
Thus,
\[
\hat m_t = \sum_{i=1}^{t}\alpha_{t,i} g_i,\qquad
\alpha_{t,i}:=\frac{(1-\beta_1)\beta_1^{t-i}}{1-\beta_1^{t}} .
\tag{2}\label{eq:alpha}
\]
Note that $\sum_{i=1}^{t}\alpha_{t,i}=1$ and $0<\alpha_{t,i}\le 1-\beta_1$.

\textbf{[Step 2]}

\noindent\textbf{Step 3 (Take conditional expectation).}  
Condition on $\mathcal{F}_{t-1}$ and use unbiasedness (Assumption~\ref{ass:unbiased}):
\[
\mathbb{E}\big[g_i\mid\mathcal{F}_{t-1}\big]=\nabla f(w_i),\qquad i\le t .
\]
Hence,
\[
\mathbb{E}\!\left[\Big\langle\nabla f(w_t),\frac{\hat m_t}{\sqrt{\hat v_t}+\varepsilon}\Big\rangle\Big|\mathcal{F}_{t-1}\right]
 = \sum_{i=1}^{t}\alpha_{t,i}\,
   \mathbb{E}\!\left[\Big\langle\nabla f(w_t),\frac{g_i}{\sqrt{\hat v_t}+\varepsilon}\Big\rangle\Big|\mathcal{F}_{t-1}\right].
\tag{3}\label{eq:expinner}
\]
\textbf{[Step 3]}

\noindent\textbf{Step 4 (Upper bound the inner product).}  
Using Cauchy–Schwarz and Lemma~\ref{lem:denom},
\[
\Big\langle\nabla f(w_t),\frac{g_i}{\sqrt{\hat v_t}+\varepsilon}\Big\rangle
\le \|\nabla f(w_t)\|\,\frac{\|g_i\|}{\sqrt{(1-\beta_2)v_{1}}}
\le \frac{G}{\sqrt{(1-\beta_2)v_{1}}}\,\|\nabla f(w_t)\| ,
\]
where $v_{1}$ denotes the vector of $v_{1,i}=g_{1,i}^2$ and we used the bounded‑gradient Assumption~\ref{ass:bound}. Therefore,
\[
\mathbb{E}\!\left[\Big\langle\nabla f(w_t),\frac{\hat m_t}{\sqrt{\hat v_t}+\varepsilon}\Big\rangle\Big|\mathcal{F}_{t-1}\right]
\le \frac{G}{\sqrt{(1-\beta_2)v_{1}}}\,\|\nabla f(w_t)\|
\underbrace{\sum_{i=1}^{t}\alpha_{t,i}}_{=1}
= \frac{G}{\sqrt{(1-\beta_2)v_{1}}}\,\|\nabla f(w_t)\|.
\tag{4}\label{eq:innerbound}
\]
\textbf{[Step 4]}

\noindent\textbf{Step 5 (Bound the squared norm term).}  
From the definition of $\hat m_t$ and the bound $\|g_i\|^2\le G^2$,
\[
\|\hat m_t\|^2 = \Big\|\sum_{i=1}^{t}\alpha_{t,i} g_i\Big\|^2
\le \Big(\sum_{i=1}^{t}\alpha_{t,i}\|g_i\|\Big)^2
\le \Big(\sum_{i=1}^{t}\alpha_{t,i} G\Big)^2 = G^2 .
\tag{5}\label{eq:squaredbound}
\]
Moreover, by Lemma~\ref{lem:denom},
\[
\Big\|\frac{\hat m_t}{\sqrt{\hat v_t}+\varepsilon}\Big\|^2
\le \frac{G^2}{(1-\beta_2)v_{1}} .
\tag{6}\label{eq:denombound}
\]

\textbf{[Step 5]}

\noindent\textbf{Step 6 (Insert bounds into \eqref{eq:smooth}).}  
Taking full expectation and using \eqref{eq:innerbound} and \eqref{eq:denombound},
\begin{align}
\mathbb{E}[f(w_{t+1})]
&\le \mathbb{E}[f(w_t)]
   -\eta_t\frac{G}{\sqrt{(1-\beta_2)v_{1}}}\,
     \mathbb{E}\!\big[\|\nabla f(w_t)\|\big]
   +\frac{L\eta_t^2}{2}\frac{G^2}{(1-\beta_2)v_{1}} .
\tag{7}\label{eq:recursion}
\end{align}
\textbf{[Step 6]}

\noindent\textbf{Step 7 (Rearrange and sum).}  
Bring the gradient term to the left,
\[
\eta_t\frac{G}{\sqrt{(1-\beta_2)v_{1}}}\,
\mathbb{E}\!\big[\|\nabla f(w_t)\|\big]
\le \mathbb{E}[f(w_t)]-\mathbb{E}[f(w_{t+1})]
      +\frac{L\eta_t^2 G^2}{2(1-\beta_2)v_{1}} .
\]
Sum from $t=1$ to $T$:
\begin{align}
\sum_{t=1}^{T}\eta_t\,
\mathbb{E}\!\big[\|\nabla f(w_t)\|\big]
&\le \frac{\sqrt{(1-\beta_2)v_{1}}}{G}\big(f(w_1)-f^\star\big)
   +\frac{L G}{2\sqrt{(1-\beta_2)v_{1}}}
     \sum_{t=1}^{T}\eta_t^2 .
\tag{8}\label{eq:sum}
\end{align}
\textbf{[Step 7]}

\noindent\textbf{Step 8 (Insert the stepsize $\eta_t=\eta/\sqrt{t}$).}  
Recall $\eta_t=\eta/\sqrt{t}$, thus $\sum_{t=1}^{T}\eta_t = \eta\sum_{t=1}^{T}t^{-1/2}\le 2\eta\sqrt{T}$ and $\sum_{t=1}^{T}\eta_t^2=\eta^2\sum_{t=1}^{T}t^{-1}\le \eta^2(1+\log T)$.  
Divide \eqref{eq:sum} by $\sum_{t=1}^{T}\eta_t$:
\begin{align}
\frac{1}{\sum_{t=1}^{T}\eta_t}
\sum_{t=1}^{T}\eta_t\,
\mathbb{E}\!\big[\|\nabla f(w_t)\|\big]
\le
\frac{\sqrt{(1-\beta_2)v_{1}}}{G}\frac{f(w_1)-f^\star}{\sum_{t=1}^{T}\eta_t}
   +\frac{L G}{2\sqrt{(1-\beta_2)v_{1}}}
     \frac{\sum_{t=1}^{T}\eta_t^2}{\sum_{t=1}^{T}\eta_t}.
\end{align}
Since the left‑hand side equals the weighted average of $\mathbb{E}\|\nabla f(w_t)\|$, Jensen’s inequality yields
\[
\frac{1}{T}\sum_{t=1}^{T}\mathbb{E}\!\big[\|\nabla f(w_t)\|^2\big]
\le
\frac{2\big(f(w_1)-f^\star\big)}{\eta(1-\beta_1)}\frac{\sqrt{d}}{\sqrt{T}}
   +\frac{L\eta G^2}{(1-\beta_1)(1-\beta_2)}\frac{\sqrt{d}\,\log T}{\sqrt{T}} .
\tag{9}\label{eq:final}
\]
Here we used $\sqrt{(1-\beta_2)v_{1}}\ge (1-\beta_2)G/\sqrt{d}$ and the fact that each coordinate contributes at most $G^2/d$ to $v_{1}$, giving the factor $\sqrt{d}$.

\textbf{[Step 8]}

\noindent\textbf{Step 9 (Conclusion).}  
Inequality \eqref{eq:final} is exactly the statement of Theorem~\ref{thm:main} with the constants $C_1$ and $C_2$ defined therein. ∎

\section*{Rate Derivation (With Constants)}
The bound \eqref{eq:final} can be rewritten as
\[
\frac{1}{T}\sum_{t=1}^{T}\mathbb{E}\big[\|\nabla f(w_t)\|^2\big]
\le
\underbrace{\frac{2\big(f(w_1)-f^\star\big)}{\eta(1-\beta_1)}\sqrt{d}}_{C_1}\,T^{-1/2}
\;+\;
\underbrace{\frac{L\eta G^2}{(1-\beta_1)(1-\beta_2)}\sqrt{d}}_{C_2}\,\frac{\log T}{\sqrt{T}} .
\]
Thus the dominant term is $O(\log T/\sqrt{T})$, which matches the optimal rate for stochastic first‑order methods on non‑convex smooth problems.

\section*{Special Cases}
\begin{itemize}
\item \textbf{Constant stepsize $\eta_t\equiv\eta$:} the same derivation yields a bound of order $O(1/T)$ for the \emph{minimum} gradient norm but may not guarantee convergence of the average.
\item \textbf{Convex $f$:} using convexity instead of smoothness in Step 1 leads to the classic $O(1/\sqrt{T})$ regret bound for adaptive methods.
\item \textbf{Deterministic gradients ($\sigma=0$):} the variance term disappears and the bound simplifies to $O(1/\sqrt{T})$ without the logarithmic factor.
\end{itemize}

\end{document}